\documentclass[11pt]{article}
\usepackage{icmlko}
\begin{document}

\title{원리적인 판이 손으로 고른 것을 못 이겼다 --- 그리고 그 실패가 앞 채택을 다시 증명했다}
\author{월드모델 실험실}
\date{노트 510 · 2026년 8월 3일}
\maketitle

\begin{abstract}
노트 509 가 \texttt{TRAINCAP = \{"만화": 0.5\}} 를 채택하면서 남긴 것 ---
\emph{도메인마다 손으로 문턱을 정하는 것보다 학습 행을 유보 무게에 맞춰
가중하는 쪽이 원리적이다.} 만들어 재 봤다.

행 무게 $=$ (그 도메인 유보 몫 $\div$ 학습 몫), $[0.2, 2.0]$ 으로 자름.
\textbf{유보 행 수만 쓰고 라벨은 안 본다} --- 채점 목적함수의 정의다.
만화 $0.25$ · 세계애니 $0.61$ · 애니 $1.18$ · 웹툰 $1.31$ · 모바일 $1.53$
· 게임 $\cdot$ 도서 $\cdot$ 팝업 $\cdot$ 시장팝업 $\cdot$ 아이돌 $2.00$(잘림).

\textbf{사전 등록(``가중이 만화 절반보다 판이 더 오른다'')은 반증} ---
가중만으로는 판 $-0.0002$($10/20$, $0$ 을 문다)다.

그런데 \textbf{가중 $+$ 만화 절반}이 판은 비기면서 KR 만화를
$+0.0203$($19/20$) 사는 것처럼 보였다. 노트 426 거울 자리라 씨앗 $40$ 으로
확인했다.

\begin{center}
\begin{tabular}{lrrl}
\hline
자 & 차 & 부호 & $95\%$ 구간 \\
\hline
판 & $-0.0008$ & $14/40$ & $[-0.0013, -0.0003]$ \\
날짜 통제 판 & $-0.0011$ & $8/40$ & $[-0.0017, -0.0006]$ \\
무리 안 판 & $+0.0008$ & $27/40$ & $[+0.0001, +0.0015]$ \\
KR 만화 & $+0.0186$ & $38/40$ & $[+0.0141, +0.0231]$ \\
비게임 앱 & $-0.0020$ & $14/40$ & $[-0.0047, +0.0006]$ \\
\hline
\end{tabular}
\end{center}

씨앗 $20$ 에서 ``$0$ 을 문다''던 판과 날짜 통제가 씨앗 $40$ 에서
\textbf{둘 다 확실한 음수}가 된다. 사전 등록은 그 둘이 $0$ 을 물 것을
요구했다 --- \textbf{기각.}
\end{abstract}

\begin{figure}[t]
\centering
\includegraphics[width=\textwidth]{figs/f1.pdf}
\caption{\textbf{가} --- 씨앗 $20$ 훑기. \textbf{나} --- 씨앗 $40$ 확인.
\textbf{다} --- 행 무게.}
\end{figure}

\section{노트 426 거울로도 안 된다}

판을 $0.0008$ 잃고 집 밖을 사는 거래로 읽어도 조건 ②(집 밖 \emph{두} 짝이
확실히 움직인다)를 못 넘는다 --- KR 은 $[+0.0141, +0.0231]$ 로 확실한데
\textbf{앱은 $[-0.0047, +0.0006]$ 로 $0$ 을 문다.}

\begin{quote}
\textbf{읽으면} --- 유보 무게 가중은 \emph{한 짝만} 산다. 노트 465 가
``집 밖 갈림은 진짜 수집 단위다''라고 했듯 두 짝은 서로 다른 것을 원하고,
이 변경은 그중 하나에만 맞는다.
\end{quote}

\section{실패가 앞 채택을 다시 증명했다}

같은 실행의 셋째 팔이 ``둘 다 끔''(가중도 만화 절반도 없는 상태)이었다.

\begin{center}
\begin{tabular}{lrrrrr}
\hline
& 판 & 날짜 통제 & 무리 안 & KR 만화 & 비게임 앱 \\
\hline
둘 다 끔 & $-0.0016$ & $-0.0025$ & $-0.0016$ & $\mathbf{-0.0385}$ & $\mathbf{-0.0120}$ \\
\hline
\end{tabular}
\end{center}

\textbf{만화 절반을 빼면 자 다섯이 전부 내려간다.} 노트 509 의 채택이
\emph{다른 실험의 대조군으로서} 독립적으로 재확인됐다 --- 특히 KR
$-0.0385$ 는 노트 509 가 잰 $+0.0274$ 보다 크다.

\section{그래도 남는 것}

\textbf{무게 자체는 합리적이었다.} 만화 $0.25$ 는 손으로 고른 $0.5$ 보다
더 세게 줄이는 방향이고, 실제로 KR 을 더 산다($+0.0186$). \textbf{틀린
것은 방향이 아니라 크기와 범위다} --- 여덟 도메인 중 다섯이 $2.0$ 에서
잘렸고(팝업은 원래 $14.6$), 그 작은 도메인들을 두 배로 키운 것이 판을
깎았을 것이다.

\begin{quote}
\textbf{다음에 시험할 것} --- 무게를 \textbf{줄이는 쪽으로만} 쓴다
($\min(w, 1)$). 만화 $0.25$ · 세계애니 $0.61$ 만 적용하고 아무것도 안
키운다. \textbf{사전 등록해 둔다.}
\end{quote}

\section*{남은 것}

\begin{itemize}
\item \textbf{줄이는 쪽 가중만}($\min(w,1)$) --- 위에 사전 등록.
\item \textbf{앱이 왜 안 움직이나} --- KR 은 $38/40$ 인데 앱은 $14/40$
이다. 두 짝이 다른 것을 원한다(노트 465).
\item \textbf{TRAINCAP 절반이 최적인가} --- $0.25$ 를 아직 안 재 봤다.
가중 결과가 그쪽을 가리킨다.
\end{itemize}

\end{document}
